\documentclass[11pt,a4paper]{article}
\usepackage[margin=2.3cm]{geometry}
\usepackage{microtype}\emergencystretch=1em
\usepackage{amsmath,amssymb}
\usepackage{booktabs}
\usepackage{graphicx}
\usepackage{xcolor}
\usepackage[inline]{enumitem}
\usepackage{hyperref}
\hypersetup{colorlinks=true,linkcolor=blue!50!black,urlcolor=blue!50!black,citecolor=blue!50!black}
\setlist{itemsep=2pt,topsep=3pt}

\newcommand{\code}[1]{\texttt{#1}}
\newcommand{\cm}{\ensuremath{\mathrm{cm}^{-1}}}
\newcommand{\ucm}{\ensuremath{\mu\mathrm{cm}^{-1}}}
\newcommand{\verdict}[1]{\par\smallskip\noindent\fbox{\parbox{\dimexpr\linewidth-2\fboxsep-2\fboxrule}{\textbf{Result.} #1}}\par\smallskip}

\title{A Critical Review of MARVEL and \textsc{MARVEL\_GNN}\\[4pt]
\large Is a trained graph neural network better than Dijkstra and other deterministic rules?}
\author{Internal review (critical position)}
\date{15 October 2026}

\begin{document}
\maketitle

\begin{abstract}
\noindent
This report takes a critical position on purpose. It examines the MARVEL method
(Furtenbacher, Cs\'asz\'ar \& Tennyson 2007; Furtenbacher \& Cs\'asz\'ar 2012). It
also examines the graph neural network (GNN) that this repository adds to MARVEL.
The GNN has one shared encoder and four heads: uncertainty, outlier detection, orphan
linkage and label correction.
The review used the code at commit \code{709c4f6} and the working tree. It ran the
shipped checkpoint (\code{data/marvel\_gnn.pt}) on the same validation data, seeds and
metrics as \code{gnn/train.py}. The results agree with the project numbers to all
printed digits. The review then added simple deterministic rules as baselines. The
project did not run these baselines.
For three of the four heads, a one-line deterministic rule gives the same result as
the GNN or a better result. Each rule uses only data that the GNN already receives as
an input feature. For the uncertainty head, the project reports a gain of 2--3
orders of magnitude over legacy MARVEL. When the deterministic rule gets one global
scale factor, this gain decreases to a factor of 1.3--1.8 in NLL.
The training target of the uncertainty head cannot see the measurement uncertainty
of a non-redundant transition. Thus the GNN gives some levels an uncertainty that is
smaller than the uncertainty of the only transition to that level. This is
physically not possible.
MARVEL also has problems. Its ``bootstrap'' is not a statistical bootstrap. Its
Dijkstra uncertainty has no clear statistical meaning. Its outlier rule cannot see
errors on bridge transitions.
\end{abstract}

\tableofcontents
\bigskip

%=======================================================================
\section{Scope and method}
%=======================================================================

\paragraph{Items examined.}
The review examined these items:
\begin{itemize}
\item \code{src/marvel\_gnn/core/}: parser, components, WLS solver and legacy
  uncertainty.
\item \code{src/marvel\_gnn/gnn/}: features, model, corruption, label correction,
  linkage and training.
\item \code{scripts/eval\_so\_calibration.py}, \code{marvel\_gnn\_plan.md},
  \code{STATUS.md}, \code{README.md} and the training logs in \code{data/}.
\item The original C++ code (\code{Source\_Code\_CPP/MARVEL4.1.cpp}).
\end{itemize}

\paragraph{Scripts.}
This directory contains two scripts:
\begin{itemize}
\item \code{skeptic\_baselines.py} loads the shipped checkpoint. It makes the
  validation sets again with \code{train.prepare}, \code{train.prepare\_corrupted} and
  \code{linkage.prepare\_linked}. It uses the \emph{same seeds} as \code{train.py}.
  It calculates the scores of the GNN and of the deterministic rules on the same
  samples.
\item \code{seed\_run.py} trains the full model five times. Only the \code{torch}
  seed changes. All data seeds stay the same. Thus the script measures only the
  variation that comes from training.
\end{itemize}
The GNN, legacy and published-$e_E$ results agree with
\code{marvel\_gnn\_plan.md} (section of 2026-08-24). For example, the GNN NLL is
7.189/6.445 and the legacy NLL is 2123.58/51.71 on \code{12C17O}/\code{13C18O}. Thus
the comparison is correct.

\paragraph{Rule for comparison.}
A claim that ``the GNN is better than MARVEL'' needs a comparison with the
\emph{strongest} simple alternative. A comparison with the weakest alternative is not
sufficient. This review supplies the strong alternatives. Each deterministic rule
uses only data that the GNN already receives as input features. Some rules also use
one global scale factor. The review fits this factor on the \emph{training}
isotopologues only.

%=======================================================================
\section{Summary}
%=======================================================================

\begin{table}[h]
\centering\small
\resizebox{\linewidth}{!}{\begin{tabular}{@{}llccc@{}}
\toprule
Head & Metric & GNN & Simple deterministic rule & Rule \\
\midrule
Uncertainty & NLL $\downarrow$ (\code{12C17O}/\code{13C18O}) & 7.19 / 6.44 & 9.23 / 11.42 & WLS covariance $\times$ Birge $\times$ 1 scale factor \\
Outlier & AP $\uparrow$ & 0.210 / 0.167 & \textbf{0.230 / 0.269} & raw $|r|/u$ (input feature 3) \\
Label correction & top-1 $\uparrow$ & 97.1\% / 96.7\% & \textbf{97.1\% / 97.2\%} & $\arg\min |E_c - E_l - \nu|/u$ \\
Orphan linkage & MRR $\uparrow$ & 0.128 / 0.052 & \textbf{0.573 / 0.094} & $|\Delta J| = 1$ first, random order for equal scores \\
\bottomrule
\end{tabular}}
\caption{The GNN and the simplest correct deterministic rule, on the same validation
samples as \code{train.py}. Bold text shows where the rule is equal to the GNN or
better.}
\label{tab:summary}
\end{table}

\begin{enumerate}
\item \textbf{Outlier detection.} The project did not compare the GNN with a residual
  threshold. This review did the comparison. The raw normalized residual gives a
  higher average precision than the GNN on \emph{both} validation isotopologues. This
  residual is already the third edge feature of the GNN (\S\ref{sec:outlier}).
\item \textbf{Label correction.} The project calls this head its ``strongest head''.
  But one of its input features, used alone, gives the same result
  (\S\ref{sec:labelfix}).
\item \textbf{Orphan linkage.} A ranking by the selection rule $\Delta J = \pm1$ is
  better than the GNN. The MRR is 4.5 times higher on \code{12C17O} and 1.8 times
  higher on \code{13C18O} (\S\ref{sec:linkage}).
\item \textbf{Uncertainty.} The project scores the legacy uncertainty against the
  GNN training target. The legacy uncertainty gets no scale factor. One fitted scale
  factor removes almost all of the difference. Also, the target does not include the
  uncertainty of non-redundant transitions. On \code{13C18O}, the GNN uncertainty of
  each of the 35 leaf levels is smaller than the uncertainty of its only transition
  (\S\ref{sec:unc}).
\item \textbf{Non-determinism.} The GNN inference is almost deterministic. On the GPU,
  the variation between runs is approximately $10^{-6}$ in $\log\sigma$. The real
  problem is different. The trained function changes with the seed, the data split and
  the stop epoch. Also, the GNN has no audit trail, no guaranteed physical properties
  and no estimate of its own (epistemic) uncertainty (\S\ref{sec:nondet}).
\item \textbf{MARVEL.} The MARVEL uncertainty method (Dijkstra and ``bootstrap'') is a
  heuristic. An exact, fast, deterministic method is available: the WLS covariance,
  which is the effective resistance of the network. Neither this project nor the C++
  code uses it (\S\ref{sec:marvel}).
\end{enumerate}

%=======================================================================
\section{Problems in MARVEL}\label{sec:marvel}
%=======================================================================

\subsection{The model is linear and Gaussian, thus an exact uncertainty is available}
MARVEL solves a weighted least-squares (WLS) problem. The incidence matrix is $B$.
Each row of $B$ is a transition and each column is a level. The weights are
$w_e = 1/u_e^2$:
\[
\hat E = \arg\min_E \sum_e w_e\,(\nu_e - (E_{u(e)} - E_{l(e)}))^2,
\qquad L = B^\top W B .
\]
The matrix $L$ is a weighted graph Laplacian. MARVEL assumes independent Gaussian
errors with standard deviation $u_e$. With this assumption, the covariance of
$\hat E_i - \hat E_0$ has an exact closed form:
\[
\operatorname{Var}(\hat E_i - \hat E_0) = R_{\mathrm{eff}}(i,0)
  = (e_i - e_0)^\top L^{+} (e_i - e_0).
\]
This value is the \emph{effective resistance} between level $i$ and the ground level.
The electrical network has a conductance $w_e$ on each edge. The function
\code{gnn/data.py:edge\_leverages} already calculates $L^{+}$ for the leverage
feature. Thus the project calculates the exact uncertainty but does not use it.

\subsection{The Dijkstra uncertainty is a loose upper limit}
Legacy MARVEL gives $\sigma_i^{\mathrm{SP}} = \sqrt{\min_{\text{paths}} \sum u_e^2}$.
This value is the resistance of the single best path. Rayleigh monotonicity gives
$R_{\mathrm{eff}} \le R_{\mathrm{shortest\ path}}$. The two values are equal only
when the network is a tree. Thus the Dijkstra method has these problems:
\begin{itemize}
\item \textbf{It ignores redundancy.} Parallel paths decrease the true variance. They
  do not decrease the Dijkstra value. The best-determined levels get the largest
  overestimate.
\item \textbf{It ignores inconsistency.} The method uses only $u_e$. It does not use
  residuals. A level with contradictory transitions gets the same $\sigma$ as a level
  with fully consistent transitions.
\item \textbf{It depends on the root.} The result changes when the reference level
  changes. For a floating component, the reference level is arbitrary.
\end{itemize}
Thus the question ``why a GNN and not Dijkstra?'' starts from an incorrect premise.
The correct deterministic comparison is not Dijkstra. It is $R_{\mathrm{eff}}$, with
a goodness-of-fit scale factor.

\subsection{The ``bootstrap'' is not a statistical bootstrap}
The code at \code{MARVEL4.1.cpp:1095--1110} changes each uncertainty to
$u_e \leftarrow k\,u_e^{\mathrm{orig}}$, with $k \sim \mathcal U\{1,\dots,10\}$. Then it
solves again. It gives $\sqrt2\max(|E-\tilde E|, s)$ as the result.
\begin{itemize}
\item It does not resample the \emph{data}. It changes only the \emph{weights}. If
  the data are fully consistent, each new solution gives the same energies. Then the
  ``bootstrap'' uncertainty is zero, for all values of the measurement uncertainties.
  The method measures the sensitivity to the weights. This is a sign of
  inconsistency, not a sampling distribution.
\item There is no derivation for the range $1..10$, the factor $\sqrt2$ or the
  maximum of the two values.
\item \textbf{The results are not reproducible.} The code initializes the random
  generator from \code{std::random\_device} (\code{mt19937 gen(rd())}). There is no
  option to set the seed. Thus legacy MARVEL is \emph{also} non-deterministic. If
  non-determinism is a problem for the GNN, it is also a problem for the published
  MARVEL uncertainties.
\end{itemize}

\subsection{The combined uncertainty counts the same uncertainty many times}
The function \code{combined\_unc} runs a second Dijkstra search. In this search, the
cost to \emph{enter} a level is $\max(\sigma^{\mathrm{SP}}, \sigma^{\mathrm{boot}})^2$.
Each level cost already includes the uncertainty of the path to that level. Thus the
sum along a path adds the same upstream uncertainty many times. The result has no
statistical meaning. It increases approximately with the square of the path length.

\subsection{The outlier rule cannot see leverage}
MARVEL puts each transition into one of these groups: WRONG, VERY BAD,
VERY BAD\_100 or VERY BAD\_1000. It uses a fixed ratio $|r_e|/u_e$
(\code{MARVEL4.1.cpp:64--90, 730}). But the fit absorbs a fraction $h_e$ of each
error. This fraction is the leverage. Thus the residual is $(1-h_e)$ times the true
error. A bridge has $h_e = 1$, so its residual is \emph{always} zero. The project
measured that 44--52\% of the transitions in the small validation networks are
bridges.
Standard regression diagnostics are available and deterministic. Examples are the
studentized deleted residual, Cook's distance and a ``cannot verify'' flag for
$h \approx 1$. MARVEL uses none of them. In the usual MARVEL workflow, a person
increases the uncertainties until the flags stop. This can hide an error behind other
errors (masking). It can also flag correct transitions (swamping).

\subsection{Adjusted uncertainties make the published data consistent}
Published MARVEL input files contain \emph{adjusted} uncertainties. A person increases
the uncertainties until the network is consistent. Thus a method that trains on the
residuals of these data learns from data that a person already made consistent. This
is a circular process.
Also, the deposited CO data are not all consistent. This review calculated the Birge
ratio $\sqrt{\chi^2/\nu}$ of each main component:
0.72 (\code{13C16O}), 0.81 (\code{12C18O}), 0.28 (\code{13C18O}), 4.9 (\code{12C17O})
and \textbf{30.9} (\code{13C17O}). For \code{13C17O}, the residuals are 31 times the
stated uncertainties. Thus this training isotopologue is not ``clean'' in a
statistical sense. This component also sets the uncertainty scale of the GNN
(\S\ref{sec:unc}).

\subsection{Structural limits}
\begin{itemize}
\item \textbf{The label is the identity.} The label string identifies the level. An
  incorrect label does not cause an inconsistency. It makes a different node. This
  node is frequently a new leaf with a zero residual. In regions with resonance mixing
  (for example the CO\textsubscript{2} Fermi index $r$), a label is not physically
  unique.
\item \textbf{Floating components.} The default is \code{--minsize 200}. With this
  value, MARVEL removes small components and gives no message. With a small
  \code{minsize}, each component gets its own arbitrary zero. MARVEL cannot connect
  these zeros.
\item \textbf{Selection-rule limits.} In SO, the e/f parity split (2026-09-14) is a
  selection-rule limit. No electric-dipole transition can connect the two parities.
  Thus no method, classical or learned, can connect these two networks with
  transition data.
\item \textbf{Systematic errors.} In real spectroscopy, the largest errors usually come
  from each source as a whole. Examples are calibration offsets, pressure shifts and
  unit errors. These errors are not independent. Thus WLS with one weight for each
  transition gives uncertainties that are too small.
\end{itemize}

\subsection{People found the real errors in this project, not models}
Each real data error in this project was found by residual calculations, physical
checks or code review. A learned model found none of them. The real errors were:
\begin{itemize}
\item SO sources in MHz with the label \cm\ (2026-09-14).
\item An incorrect F-label formula for 22HeStLy (2026-09-18).
\item MHz data in the exported \code{MARVEL\_trans.txt} files.
\item Exchanged CO\textsubscript{2} deposit files.
\item The uninitialized variable \code{hasSegmentfile} in the C++ code.
\end{itemize}
These failure types are the important types in practice. The synthetic corruption
model of the GNN contains none of them (\S\ref{sec:outlier}).

%=======================================================================
\section{Is a GNN necessary?}\label{sec:why}
%=======================================================================

The forward problem is linear and the estimator is exact. A model can learn only the
parts where the WLS assumptions are wrong. These parts are: incorrect uncertainties,
incorrect transitions and incorrect labels. For each part, a known deterministic
method is available. Each method is easy to interpret.

\begin{center}\small
\begin{tabular}{@{}lp{0.55\linewidth}l@{}}
\toprule
Task & Deterministic method & Does the GNN add value? \\
\midrule
Uncertainty & $L^{+}$ covariance; Birge or per-source scale factors; hierarchical variance model & not shown (\S\ref{sec:unc}) \\
Outliers & Studentized residual, leverage, robust (Huber/IRLS) fit & no (\S\ref{sec:outlier}) \\
Label correction & Combination-difference test on the candidate labels & no (\S\ref{sec:labelfix}) \\
Orphan linkage & Effective-Hamiltonian or Dunham prediction of absolute energies & no (\S\ref{sec:linkage}) \\
\bottomrule
\end{tabular}
\end{center}

A GNN can approximate many functions of local graph neighbourhoods. But here, the
correct function is either known in closed form (WLS covariance, residual
diagnostics), or it needs data that the GNN cannot get (absolute energies across
components). Thus the GNN must show that it adds value. The data in this report show
that it does not.

%=======================================================================
\section{Non-determinism: GNN and Dijkstra}\label{sec:nondet}
%=======================================================================

A common objection is: ``a GNN is non-deterministic, but Dijkstra is deterministic.''
This objection needs a precise statement. In its simple form, it is easy to reject.

\subsection{GNN inference is almost deterministic}
On the CPU, two forward passes of the shipped checkpoint give identical results. On
the GPU (RTX 2070, with the CUDA build of the project), the review did 20 forward
passes on \code{13C16O}. The maximum difference was
$\max|\Delta\log\sigma| = 1.9\times10^{-6}$. This difference comes from atomic
scatter-add operations. The difference between GPU and CPU was $2.4\times10^{-6}$.
These values give a relative change in $\sigma$ of approximately $10^{-6}$. This change
is not important. Also, the legacy MARVEL bootstrap is less reproducible than this.

\subsection{The data do not fully determine the trained function}
The real problem is this: $\sigma_i$ depends on the training run, not only on the data.
\begin{itemize}
\item \textbf{Seed.} Table~\ref{tab:seeds} shows five training runs. The data, the
  split and the corruption seeds are the same in each run. Only the seed for the
  initial weights changes.
\item \textbf{Data split.} The project moved one isotopologue between the training set
  and the validation set (\code{13C17O}$\leftrightarrow$\code{13C18O}).
  \code{12C17O} was in the validation set in \emph{both} runs. Its label-correction
  accuracy changed from 82.9\% to 97.1\%. The ratio between the legacy NLL and the GNN
  NLL changed from approximately 600 to approximately 8. Thus the results describe the
  split as much as they describe the method.
\item \textbf{Stop epoch.} The validation NLL decreases to 90 at epoch 300. Then it
  increases to 265 at epoch 950 (\code{data/train\_v6.log}). Thus the model overfits
  strongly. The project selects the checkpoint with the \emph{same} isotopologues that
  it then reports as validation results. Thus there is no independent test set.
\item \textbf{Version changes.} Each new training run changes all reported
  uncertainties. MARVEL data go into reference databases such as ExoMol and HITRAN. To
  reproduce one $\sigma$, a database must keep the weights, code, seeds and hardware
  details for each version.
\end{itemize}

% placeholder
\begin{table}[h]\centering\small (seed results pending)\caption{Seed variance.}\label{tab:seeds}\end{table}


\subsection{Properties that Dijkstra and WLS guarantee, but the GNN does not}
\begin{enumerate}
\item \textbf{Audit trail.} Dijkstra can explain a value of $\sigma$ with a path.
  $L^{+}$ can explain it with a sum of contributions. The GNN gives only a number from
  three attention layers.
\item \textbf{Monotonicity.} In WLS, a new measurement cannot increase the formal
  variance of a level (Rayleigh). The GNN does not enforce this property. If a person
  adds a transition, the GNN $\sigma$ can increase.
\item \textbf{Lower limit.} If one transition connects a level, then
  $\sigma \ge u_e$. The GNN does not obey this limit (\S\ref{sec:leaf}).
\item \textbf{Scale covariance.} If all $u_e$ are multiplied by $c$, each
  deterministic $\sigma$ is also multiplied by $c$. The GNN receives
  $\log_{10}u/10$ through fixed normalizations and a trained bias. On SO, the GNN gives
  an almost constant $\sigma \approx 5\times10^{-4}$ \cm\ for all networks (note of
  2026-09-17).
\item \textbf{Epistemic uncertainty.} The project uses no ensemble and no Bayesian
  method. Thus the GNN gives $\sigma$ without an uncertainty on $\sigma$. This is also
  true outside the training distribution (SO, CO\textsubscript{2}), where the error is
  largest.
\end{enumerate}

\verdict{``Non-determinism'' is not the correct word for the problem. Also, legacy
MARVEL is non-deterministic. The real problem is that the data do not fully determine
the GNN output: the seed, the split and the epoch all change it. Also, the GNN has no
audit trail. It does not obey properties that the deterministic estimators obey by
design.}

%=======================================================================
\section{Uncertainty head}\label{sec:unc}
%=======================================================================

\subsection{The target measures inconsistency, not uncertainty}
The training target comes from \code{data.py:refit\_error\_matrix}. The procedure
removes 15\% of the transitions and solves again. Then it records
$E_i^{\mathrm{masked}} - E_i^{\mathrm{full}}$. This procedure has three structural
problems:
\begin{enumerate}
\item \textbf{Consistent data give a zero target.} If all transitions agree exactly
  with the solution, each new solution gives the same energies. Then the target
  $\sigma$ is zero, for all values of $u_e$. Thus the head learns how inconsistent the
  data are. It does not learn how well the data determine each level.
\item \textbf{The scale is not defined.} The procedure is a delete-$d$ jackknife
  without the $(n-d)/d$ correction. For a correct model and a redundant network,
  $\sigma_{\mathrm{shift}} \approx \sqrt{d/(n-d)}\,\sigma \approx 0.42\,\sigma$. The plan
  says that the target is ``sensitivity under 15\% removal, not an absolute error
  scale''. But the project then compares the head with the published $e_E$, which has
  an absolute scale.
\item \textbf{Censoring.} Sometimes the procedure removes the only transition to a
  level. Then the level is not in the ground component, and its value is NaN. Thus the
  loss does not include the samples where the level is least certain.
\end{enumerate}

\subsection{The target of a leaf level does not include its own transition}\label{sec:leaf}
Let $\ell$ be a level with only one transition $e$, to parent level $p$. When $e$
stays in the sample, the solution is $\hat E_\ell = \hat E_p + \nu_e$. Thus the target
for $\ell$ is \emph{identical} to the target for $p$. The value $u_e$ is not in the
target. The physically correct uncertainty is $\sqrt{\sigma_p^2 + u_e^2}$. The table
shows the values for the validation data:

\begin{center}\small
\begin{tabular}{@{}lcc@{}}
\toprule
 & \code{12C17O} & \code{13C18O} \\
\midrule
levels / leaf levels & 33 / 4 & 345 / 35 \\
max $|\text{target}_\ell - \text{target}_p|$ (\ucm) & $1.9\times10^{-6}$ & $6.1\times10^{-5}$ \\
median $\sigma^{\mathrm{GNN}}_\ell / \sigma^{\mathrm{GNN}}_p$ & 0.998 & 1.001 \\
median $u_e / \sigma^{\mathrm{GNN}}_\ell$ & 0.18 & \textbf{4.54} \\
leaves with $\sigma^{\mathrm{GNN}}_\ell < u_e$ & 0/4 & \textbf{35/35} \\
\bottomrule
\end{tabular}
\end{center}

\noindent The target differences are only float32 noise. The GNN learned
$\sigma_\ell = \sigma_p$ exactly. On \code{13C18O}, the GNN uncertainty of \emph{each}
leaf is smaller than the stated uncertainty of the one transition that determines that
leaf. The median factor is 4.5. More training cannot correct this problem, because
the loss function rewards it.

\subsection{The comparison with legacy MARVEL is not fair}
The project scores the legacy $\sigma$ and the published $e_E$ with NLL on the GNN
target distribution. This distribution is the masked-refit shifts on the validation
isotopologues. The legacy $\sigma$ is a full-propagation uncertainty, with a very
different scale. It gets no scale factor. But the GNN learns the scale.
Table~\ref{tab:unc} adds deterministic uncertainty shapes. Each shape gets \emph{one}
global scale factor. The review fits this factor by closed-form maximum likelihood on
the three training isotopologues.

\begin{table}[h]
\centering\small
\begin{tabular}{@{}lrrrr@{}}
\toprule
& \multicolumn{2}{c}{\code{12C17O} (33 levels)} & \multicolumn{2}{c}{\code{13C18O} (345 levels)} \\
\cmidrule(lr){2-3}\cmidrule(lr){4-5}
Source of $\sigma$ & NLL & 1$\sigma$ cov. & NLL & 1$\sigma$ cov. \\
\midrule
GNN (shipped checkpoint) & \textbf{7.19} & 93.3\% & \textbf{6.44} & 97.8\% \\
legacy Dijkstra + bootstrap (project value) & 2123.58 & 85.8\% & 51.71 & 99.4\% \\
published $e_E$ & 8.33 & 96.1\% & 13.06 & 99.6\% \\
\midrule
Dijkstra only, no scale factor & 2864.34 & 83.7\% & 179.48 & 99.3\% \\
exact WLS $\sqrt{R_{\mathrm{eff}}}$, no scale factor & 2882.05 & 82.8\% & 186.01 & 99.3\% \\
exact WLS $\times$ Birge ratio & 124.18 & 87.5\% & 2247.36 & 98.5\% \\
\midrule
Dijkstra $\times$ 1 fitted scale factor & 11.11 & 100\% & 16.47 & 100\% \\
exact WLS $\times$ 1 fitted scale factor & 10.94 & 100\% & 15.78 & 100\% \\
exact WLS $\times$ Birge $\times$ 1 fitted scale factor & 9.23 & 99.6\% & 11.42 & 99.7\% \\
\bottomrule
\end{tabular}
\caption{Uncertainty NLL (a lower value is better; $\sigma$ in \ucm) on the validation
masked-refit errors from \code{train.py}. The fitted scale factors are 2782
(Dijkstra), 2790 ($R_{\mathrm{eff}}$) and 101 ($R_{\mathrm{eff}}\times$Birge).}
\label{tab:unc}
\end{table}

\paragraph{Interpretation.}
\begin{itemize}
\item One parameter decreases the Dijkstra NLL from 2864 to 11.1. Thus the gain of
  ``2--3 orders of magnitude'' comes almost fully from the scale. The project compares
  a value with no scale factor to a value with a fitted scale.
\item The real GNN gain over the best one-parameter rule is
  $\Delta\mathrm{NLL} \approx 2$ on \code{12C17O} and $\approx 5$ on \code{13C18O}. This
  gain is real but small. The review measured it on the GNN target, on two
  isotopologues. One of these isotopologues has only 33 levels.
\item The fitted scale factor is approximately 2800 times the formal WLS $\sigma$. With
  the Birge correction, it is approximately 100 times. A single maximum-likelihood
  scale also gives approximately 100\% coverage at $1\sigma$. These two results show
  that the target has heavy tails and many zeros. Most masks do not move a level. A
  small number of masks move the level through inconsistent transitions. A Gaussian
  NLL is not the correct likelihood for this distribution. Thus \emph{each} method,
  also the GNN, gives too much coverage at $1\sigma$ (93--98\%, but the target is 68\%).
\item The Spearman correlation between the GNN $\sigma$ and the exact WLS $\sigma$ is
  $-0.14$ on \code{12C17O} and $0.62$ on \code{13C18O}. On the small network, the GNN
  ranking has no relation to the formally correct ranking.
\end{itemize}

\subsection{Other problems in the evaluation}
\begin{itemize}
\item \textbf{A check that is always true.} The check
  ``$|E_{\mathrm{ours}} - E_{\mathrm{pub}}| \le \sigma$: 100\%'' is always true. The
  port gives the published energies to $10^{-8}$ \cm. Thus the left side is almost
  zero.
\item \textbf{A small, correlated validation set.} The five CO isotopologues have
  almost the same topology and QN structure. Thus ``generalization across
  isotopologues'' is interpolation. \code{12C17O} has only 33 levels and 50
  transitions.
\item \textbf{No uncertainty on the metrics.} Each result comes from one run, one split
  and one bootstrap seed. There are no confidence intervals.
\end{itemize}

\verdict{The head optimizes a target with no absolute scale. This target does not
include the measurement uncertainty of non-redundant transitions. Most of the
reported gain over legacy MARVEL comes from the scale. A one-parameter deterministic
rule gives almost the same NLL. A deterministic rule never gives
$\sigma_\ell < u_e$.}

%=======================================================================
\section{Outlier-detection head}\label{sec:outlier}
%=======================================================================

The plan does not compare the GNN with the residual-ratio rule. The plan says that
this comparison ``risks inheriting its blind spots''. But a comparison does not give
blind spots to a model. A comparison only shows if the model adds value. The review
did the comparison on the same 50 corrupted copies of each validation isotopologue:

\begin{table}[h]
\centering\small
\begin{tabular}{@{}lcccccc@{}}
\toprule
& \multicolumn{3}{c}{\code{12C17O} (base rate 3.9\%)} & \multicolumn{3}{c}{\code{13C18O} (base rate 4.9\%)} \\
\cmidrule(lr){2-4}\cmidrule(lr){5-7}
Score & AP & AP freq & AP QN & AP & AP freq & AP QN \\
\midrule
GNN & 0.210 & 0.041 & 0.295 & 0.167 & 0.027 & 0.211 \\
raw $\log(1+|r|/u)$ (edge feature 3) & \textbf{0.230} & \textbf{0.044} & \textbf{0.328} & \textbf{0.269} & 0.027 & \textbf{0.401} \\
studentized residual (edge feature 5) & 0.219 & 0.042 & 0.317 & 0.249 & 0.026 & 0.366 \\
\bottomrule
\end{tabular}
\caption{Average precision on the corrupted validation copies. One input feature is
better than the full GNN on both isotopologues. For QN changes on \code{13C18O}, it is
1.9 times better.}
\end{table}

\paragraph{Other problems.}
\begin{itemize}
\item \textbf{Training data and test data come from the same synthetic generator.}
  The generator corrupts 5\% of the transitions. Each corruption is a log-uniform
  shift of 5--500$u$ or a $J\pm1$ change on the upper level. Thus the test measures
  how well the GNN inverts its own generator. It does not measure how well the GNN
  finds real errors.
\item \textbf{The real failure types are not in the generator.} The real types are:
  unit errors (a full source scaled by $1/29979$); systematic label errors in a full
  source (the 22HeStLy F-label error affected approximately 60\% of a network, not
  5\%); exchanged deposit files; air and vacuum wavenumbers; exchanged upper and lower
  levels; incorrect e/f and $K_a/K_c$ labels; and transitions from the incorrect
  isotopologue. Each real error in this project was of one of these types.
\item \textbf{The ``structural floor'' is a deterministic fact.} No method can find an
  error on a bridge ($h=1$). A deterministic output can show $h$ for each transition.
  This needs one line of code and no model.
\end{itemize}

\verdict{The GNN is worse than one of its own input features. The decision to not
compare the GNN with the residual rule hid this result.}

%=======================================================================
\section{Label-correction head}\label{sec:labelfix}
%=======================================================================

The candidate window contains the levels with the same non-$J$ tokens and with $J$
within $\pm2$. The corruption is exactly a $J\pm1$ change on the upper level. The
evaluation uses only the ``recoverable subset'', where the true upper level is in the
window. The candidate feature \code{extra[2]} $= \log(1+|E_c - E_l - \nu|/u)$ is the
classical combination-difference test. The plan already calls it ``the dominant
feature''.

\begin{center}\small
\begin{tabular}{@{}lcccc@{}}
\toprule
& \multicolumn{2}{c}{\code{12C17O} ($n=35$, window 3.3)} & \multicolumn{2}{c}{\code{13C18O} ($n=720$, window 3.9)} \\
\cmidrule(lr){2-3}\cmidrule(lr){4-5}
Ranking method & top-1 & MRR & top-1 & MRR \\
\midrule
GNN & 97.1\% & 0.979 & 96.7\% & 0.982 \\
$\arg\min$ mismatch (one feature, no model) & \textbf{97.1\%} & \textbf{0.986} & \textbf{97.2\%} & \textbf{0.986} \\
\bottomrule
\end{tabular}
\end{center}

\begin{itemize}
\item The ``random 30\%'' baseline is too weak to be useful. The project designed the
  task so that one arithmetic check solves it.
\item The task is narrow by design. It has one QN, one direction (upper level), one
  error type ($\pm1$) and only the recoverable cases. Real label errors include
  $\Delta J = \pm2$, parity, $K_a/K_c$ exchanges, vibrational label errors in
  polyads and errors on the lower level.
\item The selection has two stages. First, the outlier head must flag the transition
  ($p > 0.5$). Then the label-correction head runs. Table~\ref{tab:summary} shows that
  the outlier head is the weak stage.
\end{itemize}

\verdict{The head gives the same result as a deterministic combination-difference
test. It adds no measurable value. It can be a useful tool, but it is not evidence
for the GNN.}

%=======================================================================
\section{Orphan-linkage head}\label{sec:linkage}
%=======================================================================

\begin{center}\small
\begin{tabular}{@{}lcccccc@{}}
\toprule
& \multicolumn{3}{c}{\code{12C17O} (109 pairs)} & \multicolumn{3}{c}{\code{13C18O} (884 pairs)} \\
\cmidrule(lr){2-4}\cmidrule(lr){5-7}
Ranking method & MRR & hits@1 & hits@5 & MRR & hits@1 & hits@5 \\
\midrule
GNN & 0.128 & 0\% & 20\% & 0.052 & 0\% & 12.2\% \\
$|\Delta J|=1$ first, random order for equal scores & \textbf{0.573} & \textbf{33\%} & \textbf{100\%} & \textbf{0.094} & \textbf{2.7\%} & \textbf{12.8\%} \\
\bottomrule
\end{tabular}
\end{center}

\begin{itemize}
\item \textbf{The GNN is worse than the selection rule that it receives as a feature.}
  The $|\Delta\mathrm{QN}|$ block is an input to the head. But a rule that uses only
  $|\Delta J|$ has an MRR 4.5 times higher than the GNN on \code{12C17O}.
\item \textbf{The task is not well posed.} The scientific result is not \emph{which}
  bridge transition the procedure removed. Without a measured frequency, each pair with
  $\Delta J=\pm1$ agrees equally well with the data. The necessary result is the
  \emph{absolute energy offset} of the orphan component. An effective-Hamiltonian or
  Dunham fit to the main component gives this offset directly. For a diatomic
  molecule, the accuracy is frequently $10^{-3}$ \cm\ or better. By design, the GNN
  never sees absolute energy across components.
\item \textbf{The training molecule has no real orphan case.} The CO networks have no
  real orphan components. The project found one real orphan case (SO). That case is a
  parity limit. No linkage model can or must connect it.
\end{itemize}

\verdict{Remove this head, or replace it with an effective-Hamiltonian prediction of
the offset. The project documentation already calls the head weak. The data show
that it is worse than a one-line rule.}

%=======================================================================
\section{Other design and engineering problems}
%=======================================================================

\begin{itemize}
\item \textbf{Multi-task coupling.} The loss weights of $10\times$ have no
  justification. Early stopping uses only the uncertainty NLL. Thus the other three
  heads use the weights from the best uncertainty epoch, which can be a bad epoch for
  them.
\item \textbf{Features.} The QN tokens are divided by 50. The code changes
  \code{e}/\code{f} to 0/1. Energies go in as $\log(1+E)/10$. The code does not
  support categorical tokens (electronic state, F1/F2/F3). For this reason, the
  \code{\_encode\_so} patch is necessary at evaluation time. The project did not test
  the ``schema-oblivious'' claim on asymmetric tops.
\item \textbf{Scalability.} \code{edge\_leverages} inverts an $n\times n$ matrix. This
  stops CO\textsubscript{2} training. The deterministic baselines need the same
  diagonal of $L^{+}$. A sparse Cholesky factorization with selected inversion
  (Takahashi recursion) calculates it exactly at CO\textsubscript{2} size. When this
  calculation is available, the deterministic baseline is also available.
\item \textbf{SO evaluation outside the training distribution.} There is no ground
  truth. The ``badness'' value (max $|r|/u$) comes from the same solution. The negative
  correlation of the legacy method has no explanation. The project already withdrew
  one earlier version of this result. A Spearman $\rho$ over hundreds of dependent
  levels, without confidence intervals, is weak evidence.
\item \textbf{Documents that are not current.} \code{README.md} gives a
  label-correction accuracy of 88.6\%/83.6\%. These values come from an old split.
  \code{STATUS.md} (2026-07-29) says that the C++ code is not built. The changes to
  the split and the path in \code{train.py} are not committed. The file
  \code{train\_v7\_13C18O.log} uses UTF-16, so standard tools cannot read it.
\end{itemize}

%=======================================================================
\section{Scientific integrity}
%=======================================================================

The project rule ``advisory only'' for linkage and label correction is correct. But
these two heads are the least important, because deterministic rules give the same
results. The uncertainty head is different. Its output replaces the published
$\sigma$, so it can go into a database without human review. This head has the
incorrect target (\S\ref{sec:leaf}). Uncertainties go into line lists (ExoMol,
HITRAN). From there, they go into calculated abundances. A $\sigma$ that is smaller
than the uncertainty of its own measurement is a scientific-integrity problem. This is
the type of problem that \code{CLAUDE.md} of this project describes.

%=======================================================================
\section{Evidence that would change the result of this review}
%=======================================================================

\begin{enumerate}
\item \textbf{Use the correct deterministic baseline.} Replace ``legacy
  Dijkstra+bootstrap'' with the exact WLS covariance. Add a Birge factor or a
  per-source variance model (a hierarchical model with one scale factor for each
  source, fitted by REML). Claim a GNN gain only when the GNN is better than this
  baseline.
\item \textbf{Correct the target.} Measure the calibration on \emph{validation
  transitions}. Predict $\nu_e$ from $\hat E_u - \hat E_l$. Test the z-score against
  $\sqrt{\sigma_u^2 + \sigma_l^2 - 2\,\mathrm{cov} + u_e^2}$. This test works at the
  measurement level. It has an absolute scale, and it includes leaf transitions.
\item \textbf{Enforce physics.} Use $\sigma_\ell \ge u_e$ as a hard lower limit. A
  better method is to predict a \emph{correction factor} on the WLS covariance. Then
  the model obeys the properties by design.
\item \textbf{Compare each head with its rule}, as in Table~\ref{tab:summary}. Report
  the GNN only where it is better.
\item \textbf{Use real errors.} Different MARVEL releases of the same molecule (for
  example the updates for H\textsubscript{2}\textsuperscript{16}O) contain transitions
  that a person removed, weighted again or labelled again. These are \emph{real}
  errors with human labels. They are a much better test set for outlier detection and
  label correction than synthetic $J\pm1$ changes. The SO errors of this project
  (units, F labels) are also real cases.
\item \textbf{Use correct statistical procedures.} Keep a test set that early
  stopping does not use. Use $k$-fold validation over isotopologues. Use 5 seeds or
  more and give the mean $\pm$ standard deviation. Use deep ensembles to estimate the
  epistemic uncertainty.
\item \textbf{Orphans.} Compare the head with an effective-Hamiltonian prediction of
  the offset (Dunham, PGOPHER, Duo). Or remove the head.
\end{enumerate}

\subsection*{Where a learned model can possibly add value}
A learned model can possibly add value in these areas, where deterministic methods
are difficult:
\begin{enumerate*}[label=(\roman*)]
\item learn the \emph{per-source} systematic-error scale from many molecules, as a
  prior that replaces manual adjustment;
\item rank suspect transitions \emph{for human review}, where only the order is
  important and the rule-based scores above are the input features;
\item transfer priors between molecules that use the same experimental method.
\end{enumerate*}
The current heads do none of these tasks. The project has not shown any of them.

%=======================================================================
\appendix
\section{How to reproduce the results}
%=======================================================================
Use a Python 3.12 environment with \code{torch}, \code{torch\_geometric},
\code{numpy}, \code{scipy} and \code{networkx}. A CPU build is sufficient. Run these
commands from the repository root:
\begin{verbatim}
python -W ignore reports/2026-10-15_skeptic_review/skeptic_baselines.py
python -W ignore reports/2026-10-15_skeptic_review/seed_run.py <seed>   # 0..4
\end{verbatim}
Both scripts use the fixed WSL path \code{/mnt/c/Code/MARVEL\_scraping/molecules/CO}
for the CO data. \code{seed\_run.py} writes its checkpoint into its own directory. It
does not write over \code{data/marvel\_gnn.pt}. The raw results are in
\code{skeptic\_results.json} and \code{seed*.log}.

\end{document}
